\subsection{Limiti}
I limiti del modello sono di tre tipi: ciò che il modello non rappresenta per scelta, ciò che rappresenta con dati incerti e ciò che rappresenta in modo non corretto.

Al primo tipo appartengono le esclusioni dichiarate fin dall'inizio. Il modello è deterministico: non ci sono guasti né perturbazioni casuali, lo stesso scenario produce sempre lo stesso esito e i risultati non hanno un intervallo di confidenza. Gli indicatori di puntualità sono per questo un limite superiore di quelli reali. Non ci sono passeggeri: il riempimento dei treni è un'ipotesi per fascia oraria, i tempi di sosta sono quelli dell'orario e non crescono con l'affollamento, e non si misura il tempo di viaggio degli utenti, che è ciò a cui un aumento di frequenza serve. Il servizio simulato è il solo servizio di Trenord: i treni merci, quelli a lunga percorrenza e l'alta velocità, che sulla stessa infrastruttura occupano capacità, non ci sono. Non sono rappresentati i passaggi a livello, e le variazioni temporanee dell'orario per lavori sono escluse.

Al secondo tipo appartengono i dati stimati. La capienza delle tratte a doppio binario deriva da una lunghezza delle sezioni di blocco uguale per tutte le linee, non rilevata. Sono stime il tempo di inversione ai capolinea, la soglia oltre la quale un treno va al ricovero e la capienza di alcuni ricoveri, fra cui quello di Milano Cadorna, reso volutamente non vincolante. Sono stimate le accelerazioni di una parte dei treni e la decelerazione di tutti. I giri treno del modello non sono i turni dell'operatore, che non sono pubblici, e il tipo di treno di ogni linea è una ripartizione in percentuale: il numero dei treni per tipo ne dipende. I costi unitari sono stime di ordine di grandezza, e per tre di essi manca una fonte verificata. Nel modello energetico sono ipotesi i servizi ausiliari, il rendimento del motore diesel e il riempimento dei treni; pendenze e curve sono trascurate, e il recupero in frenata è valutato in linea d'aria, senza lo schema elettrico della rete.

Al terzo tipo appartengono i difetti noti. Il profilo di velocità usato porta i treni sempre alla velocità massima, con accelerazione e frenata massime: i treni simulati sono più veloci dell'orario su quasi metà delle tratte e consumano più del reale, cosa che per il gasolio è stata corretta con una taratura e per l'energia elettrica no. Una corsa che non ferma in una stazione intermedia percorre, in una parte dei casi, un collegamento diretto distinto da quello dei treni che vi fermano: è un binario che non esiste, e gonfia la capacità delle linee su cui accade. Nelle stazioni di transito più cariche restano 131 assegnazioni di binario forzate, con due treni ravvicinati sullo stesso binario. Soprattutto, quando la domanda supera la capacità il modello non degrada come la rete reale: manca chi regola la circolazione, trattiene i treni a monte e sopprime le corse, e i treni si bloccano a vicenda in modo permanente. Uno scenario oltre la soglia dice che la soglia è stata superata e dove, ma non quanto peggiorerebbe il servizio.

\subsection{Sviluppi futuri}
Gli sviluppi più utili sono quelli che rimuovono un limite del terzo tipo. Il primo è una regolazione della circolazione che riproduca le decisioni di chi la governa: precedenze fra treni di categoria diversa, trattenimento a monte di un punto saturo, soppressione di una corsa. railsim prevede a questo scopo un punto di estensione, quello che decide a chi concedere un tratto conteso, che nel progetto è rimasto quello predefinito, senza priorità. Sarebbe insieme la correzione del comportamento oltre la soglia e la leva sperimentale più interessante, perché permetterebbe di confrontare politiche di precedenza diverse a parità di orario. Il secondo è l'uso del profilo di velocità adattivo già offerto da railsim, che regola la velocità sull'orario di arrivo: renderebbe più realistici i tempi di percorrenza e i consumi, e la taratura del gasolio non servirebbe più.

Altri sviluppi estendono ciò che il modello rappresenta. Con perturbazioni casuali, cioè guasti e ritardi alla partenza estratti da distribuzioni, si potrebbero ripetere le simulazioni e ottenere indicatori con una dispersione, confrontabili con la puntualità reale per tutte le cause. Con i passeggeri, che MATSim sa simulare, si misurerebbe il beneficio di uno scenario in tempo di viaggio e di attesa, e non solo il suo costo. Con i treni delle altre imprese la capacità residua dei nodi sarebbe quella vera.

Un'analisi rimasta aperta riguarda il binario unico. I risultati indicano che è il primo limite che si incontra aumentando le frequenze, e il modello permette di valutare l'effetto di un raddoppio selettivo: basterebbe modificare il numero dei binari di una tratta nei dati e ripetere la simulazione, confrontando il beneficio con il costo dell'opera.

\subsection{Conclusioni}
Il progetto ha costruito un modello della circolazione dei treni regionali e suburbani della Lombardia, a partire da dati pubblici, e un'applicazione che permette di simulare una giornata di servizio, modificarne l'orario e misurarne puntualità, energia, costi e treni necessari. Il modello fa circolare l'intero orario di un giorno feriale, 2349 corse effettuate da 343 treni, senza che alcun treno resti bloccato, e per le grandezze che rappresenta sta vicino ai dati pubblicati dall'operatore.

Alla domanda di ricerca, cioè di quanto si possa aumentare la frequenza prima che i ritardi si propaghino e quali punti fissino il limite, i risultati danno una risposta articolata. Sui tratti urbani di Milano l'attesa massima fra due treni può essere ridotta a 10 minuti nelle ore di punta, aggiungendo corse limitate a quei tratti, senza che la rete si blocchi e senza superare sulla carta la flotta disponibile. La cadenza delle singole linee sull'intero percorso resta quella attuale, tranne che per una linea. Il primo limite non è dove ci si aspetterebbe, nel tunnel del Passante, ma nel tratto a binario unico di una delle linee potenziate. Oltre queste frequenze il limite diventa il tunnel con i suoi capolinea, da cui la coda si propaga al resto della rete.

Il lavoro ha mostrato anche le proprietà che fanno della rete ferroviaria un sistema complesso. Il comportamento d'insieme non era ricavabile dalle regole dei singoli treni: ogni regola aggiunta al modello è nata da una simulazione in cui la rete si fermava per una causa non prevista, e più di una volta una regola locale, come l'ordine con cui si assegnano i binari di una stazione, ha cambiato l'esito dell'intera giornata. Gli effetti non sono proporzionali alle cause, e un errore in un solo punto della rete è rimasto invisibile finché il carico non è aumentato.

Resta il limite principale: il modello dice dove e quando la rete smette di reggere, ma non descrive in modo realistico ciò che accade dopo. Per farlo servirebbe rappresentare chi, nella rete reale, regola la circolazione; è lo sviluppo che renderebbe il modello utile anche oltre la soglia.
